\section{Meta evaluación y generalización a largo plazo}
\subsection{Detalles del entrenamiento del Meta-Judge}
El meta judge surge porque los modelos self-rewarding llegan a un plateau después de few iterations: Weston lo describe de manera gráfica: ``self-rewarding models that I showed earlier only had the top part of the pipeline; introducing meta judge builds the bottom part, lifting the plateau'' (SRT 2756-2789). La entrada del meta judge es el par draft y verified response, y la salida indica cuál de los dos vale la pena conservar; el objetivo de entrenamiento ya no depende de preferencias humanas, sino que se refuerza según el win rate en benchmarks como Alpaca Eval 2.

\begin{figure}[H]
\centering
\includegraphics[width=0.77\textwidth]{frame-04-meta.jpg}
\caption{El capítulo del Meta Judge muestra la división de trabajo en tres pasos: draft, judge y verified.\protect\footnotemark}
\end{figure}
\footnotetext{Intervalo de tiempo en el video: 01:01:30--01:01:32.}

\subsection{Datos del Meta-Judge y SelfT Evaluators}
Weston menciona que ``SelfT evaluators son los que hicimos el año pasado en el paper de SelfT''\footnote{Traducción del hablante; en el original en inglés: ``SelfT evaluators''.}, un conjunto de datos que estructura el proceso de decisión sobre la verified response: incluye tres etapas, draft, verified y judge, y la etiqueta del judge proviene de tareas de razonamiento de escritorio y no de preferencias puramente humanas (SRT 2902-2904). El judge entrenado con los SelfT evaluators logra distinguir con más precisión las respuestas que ``realmente mejoran la calidad del reasoning'', con lo cual se eleva el win rate.

\begin{knowledgebox}{Lógica de diseño de los SelfT Evaluators}
Los SelfT Evaluators se entrenan con borradores que tienen chain-of-thought y respuestas ya verificadas; el judge solo tiene que elegir la más confiable entre ambas. Un evaluator entrenado así logra, con poca supervisión humana o con datos autosupervisados, una capacidad de discriminación mayor que la de los modelos de reward tradicionales.
\end{knowledgebox}

\subsection{Interpretabilidad del razonamiento y el reto computacional}
Weston sostiene que ``transformer itself does reasoning by all these steps'', pero al mismo tiempo reconoce que la red en sí no es transparente; por eso es necesario recurrir al meta judge o a un reward verificable para rastrear el reasoning trace (SRT 3152-3204). También insiste en que el ``inference time compute for evaluation'' es el cuello de botella del self-improvement, razón por la cual el número de verified response no puede aumentarse a la ligera, pues de lo contrario el costo de inferencia se duplicaría (SRT 3236-3252).

\subsection{Resumen de la sección}
El meta judge y los SelfT evaluators dotan a los modelos self-rewarding de una capacidad de discriminación genuina: al comparar el draft con la verified response, se asegura que la señal de reward no se confunda con hallucination, y al mismo tiempo se nos recuerda que el cómputo de inferencia debe mantenerse bajo control.

